\section{合卷：晋阳城下，春秋最后一个大国被分开}

前403年，周威烈王册命韩、赵、魏为诸侯。传世材料留下的是这项册命及其年代，而不是一场可以复原的朝廷仪式；但这份名分仍有分量。它使三个本在晋国名义下行动的政治集团，能够以各自的名号同列国交往。只是册命所承认的，并不是刚刚出现的三块空地，而是已经各有军政资源、彼此难以再由晋君节制的现实。由\eventref{WQ-0403-THREEJIN}{战国·三家分晋}回望，春秋与战国之间的门槛并不在某一天忽然落下，而在于旧有的封国、家门与王室名分，先后失去了原来的扣合方式。

\subsection{晋阳之后：谁能把晋的资源继续调动起来}

晋的困境并非从智氏一役才开始。对齐、楚的长期竞争曾要求扩大动员，\eventref{SA-0588-JIN-SIX-ARMIES}{春秋·晋扩编六军}把更多军政位置分置于卿族，增强了对外作战的能力，也让统军的职位、随从和地方关系有了在家门中延续的条件。此处不宜把军制的名称或规模简单等同于后来的常备官僚制；能够稳妥把握的是，国君要借卿族组织军队，卿族也因而更深地进入军政资源的分配。

\eventref{SA-0546-MIBING}{春秋·弭兵之会}后，晋、楚之间的直接大规模冲突暂缓。它不是分晋的单一原因，却是一个使能条件：共同作战的压力减轻之后，原来为对外战争而配置的资源，更容易成为国内各家反复争夺的对象。到\eventref{SA-0497-JIN-FACTION}{春秋·晋卿族内战}，范氏、中行氏与赵氏的冲突已经表明，晋的政治问题不只是君臣间谁更受宠，而是谁能据有城邑、部众和军政职位，谁又能在冲突中取得其他家门的支持。

前453年的\eventref{WQ-0453-ZHI}{战国·三家灭智}，把这种积累推到无法调和的地步。传世叙事把危机集中在晋阳，将智伯、赵氏以及韩、魏的转向写得格外鲜明；可确证的结果是韩、赵、魏合力消灭智氏，晋国的实际控制因而分裂。把这一结局归结为智伯的性格，或把三家合作写成早已注定的联盟，都会遮蔽当时的选择：韩、赵、魏确实选择了共同对付智氏，但各人怎样权衡风险、每一次谋议怎样发生，现存材料不能逐场复原。结构性的压力在于卿族长期掌握可征发、可赏赐、可统兵的资源；触发局势的却是家门之间的冲突与结盟。

因此，周廷的承认不是分裂的起点。册命前，三家已经能够分别支配晋国旧有资源；册命后，它们获得了可以公开使用的诸侯名分。晋君的名号也没有立即消失，直到\eventref{WQ-0376-JIN-PARTITION}{战国·韩赵魏废晋静公}，韩、赵、魏才废晋静公并分其余地。前403年至前376年的间隔尤其提醒人们：名分、实际控制和旧国宗祀的终结可以不同步，不能用一纸册命抹平这一段仍在重组的政治时间。

\subsection{齐与鲁：同一压力，不同的国家形状}

齐提供了另一条通向战国的路径。前481年，\eventref{SA-0481-LIN}{春秋·获麟与陈氏政变}把陈恒杀齐简公与《春秋》经文的终篇并置在同一年。它并不证明田氏从此立即拥有齐国，也不证明数十年的演变只是一次预谋的夺国；它只使一个早已积累的事实格外清楚：齐君的存废已经可能由强大的家门决定。至\eventref{WQ-0386-TIANQI}{战国·田氏代齐}，田氏获得周廷册命，才以齐侯的身份接过齐国的对外名号与责任。

齐与晋的相似处，是资源的归属先发生变化，王室承认后到；不同处则同样重要。田氏是以一个家门取代姜齐，而韩、赵、魏是把晋的资源拆分为三国。前者没有使“齐”这个政治名称消失，后者却使“晋”逐步失去作为一个国家的实际内容。把两件事都称为“篡夺”或都称为“新国家诞生”，都太快了；它们共同提出的只是一个更具体的问题：当城邑、兵众、征发与执政职位不再由旧君室有效统合时，谁还能以一国的名义处理这些事务。

鲁的三桓使这个问题更容易看见，却也显示它并非必然通向易姓或分国。\eventref{SA-0517-LU-ZHAO-EXILE}{春秋·鲁昭公出奔}之后，鲁昭公未能在国内重建立足处；\eventref{SA-0502-YANGHU}{春秋·阳虎出奔}又显示，三桓的家臣能够借主家的邑兵与职位越级参与政治。权力不是从鲁君直线移给三桓，而是在公室、卿族和家臣的多层关系中移动。鲁国仍保有君主和国号，并不意味着公室已能直接穿透这些层次。

前498年的\eventref{SA-0498-DUO-SANDU}{春秋·堕三都}尤其能看出行政条件的边界。郈、费被堕而成未克，既说明鲁君与部分卿大夫曾试图拆除支族、家臣据兵征赋的据点，也说明一次行动不足以收回三桓长期占有的邑兵与人力。前483年的\eventref{SA-0483-LU-TIANFU}{春秋·鲁用田赋}又把田亩同军赋征调重新联结起来；这提示战事与城防需要更稳定的负担安排，却不能据此想象鲁国已经具有后世那种完整的户籍、税率或县级官僚档案。三桓下的鲁没有成为第二个晋：它留下的是国号未改而公室空心化的局面。

\begin{causalchain}[从家门竞争到列国竞争]
\term{结构性压力}是春秋国家须借封邑、卿族与家臣来统兵、守城、征发和办理外交；这些安排提高了动员能力，也让职位与资源可能在家门内沉积。\term{行动者的选择}则发生在具体危机中：韩、赵、魏合灭智氏，陈恒越过齐君，鲁君及部分卿大夫试图削弱三都；它们不能被倒写成同一套动机或同一条必经道路。

\term{使能条件}是周王室仍能提供诸侯名分，却已无力把分散的军政财赋重新收归旧国。威烈王对三晋的册命和周廷对田氏的承认，赋予既成政治体可辨认的外交与礼制位置，并未替它们建立征发、任官或统军的能力。\term{中期影响}是竞争的单位逐渐不再只是主持会盟的霸主，而是能够更稳定地组织土地、人力、军队和行政职位的国家；这正是战国各国接下来要以不同办法解决的难题。
\end{causalchain}

\begin{turningpoint}[三家分晋为什么是春秋和战国之间的门槛]
三家分晋的重要性，不在于它替整个时代划出一条毫无过渡的界线，而在于春秋最强的北方大国已无法维持一个共同的军政中心。会盟、宗法和王命仍能被援引，却不能自动决定资源归谁调度。齐的田氏接国、鲁的三桓制约公室和晋的三家分立，是同一压力下的不同结果；战国国家此后会更迫切地尝试把征发、军队和官职接到更可持续的政治中心上，也会让这种重组的成本更直接地落到城邑与家庭。
\end{turningpoint}

\begin{sourcenote}[晋阳、册命与家门上升的材料边界]
晋卿族冲突、鲁昭公出奔、阳虎出奔、堕三都和田赋，可先以《春秋》的年次为骨架，再参照《左传》与《史记》的叙事；经文能够确认的事实十分简略，不能把传世说辞当作逐字记录。三家灭智、周威烈王册命三晋的年代与结果，主要据《资治通鉴》卷一、《史记》各国世家及《六国年表》；田氏获册命则可参照《史记·田敬仲完世家》与《资治通鉴》。这些都是后出整理或后世编年，不等于现场档案。

所以，晋阳围城中的劝说、人物品评、水攻细节，以及田氏、三桓各自如何逐年经营人众和城邑，仍有材料层次与解释空间。本文只据可核的事件链讨论资源、名分与行政关系，不将戏剧性最强的版本写成唯一现场，也不把后来的战国制度倒灌到春秋国家。
\end{sourcenote}
